% =============================================================================
% Chapter 5: Preliminary Results and Discussion -- Part I
% =============================================================================
\mychapter{5. Preliminary Results and Discussion}

This chapter presents the preliminary results obtained from the preprocessing pipeline validation,
statistical analysis, parameter sensitivity analysis, and initial detection model training on the
NIST CoCr \texttt{set1sample2raw} dataset. All quantitative results reported here are measured from
the implemented pipeline as described in Chapter~3.

\mysection{5.1 Key Insights from Literature Review}
The systematic literature review produced several findings with direct implications for the
experimental design of this thesis:
\begin{itemize}
\item U-Net-based deep learning consistently outperforms classical threshold-based detection methods
  by 20--40 Dice percentage points on AM XCT data, with the largest gains observed for small defects
  and artefact-affected regions~\citep{Wong2022}.
\item 2D slice-based detection achieves within 2--4\% of full 3D volumetric performance on isotropic
  defect types such as gas porosity, at 10--50$\times$ lower computational cost~\citep{Wong2021}.
  This finding directly motivates the 2D-first approach adopted in this thesis.
\item Combined Dice-Focal loss achieves superior detection performance over either component alone
  under the class imbalance conditions typical of AM XCT data~\citep{Bellens2021}. Preliminary
  training results on NIST CoCr data confirm this finding (Section~5.5).
\item Increasing architecture complexity yields diminishing returns unless training data diversity is
  simultaneously increased~\citep{Zubayer2023}. This supports the data-centric approach adopted in
  this work, prioritising preprocessing quality and augmentation diversity over architectural
  novelty.
\item Preprocessing, particularly beam hardening correction and percentile normalisation,
  substantially affects detection performance and is frequently under-reported in published studies.
  This thesis addresses this gap with a comprehensive parameter sensitivity analysis
  (Section~5.2).
\item Training label quality is the dominant performance factor, exceeding the contribution of
  architecture choice~\citep{Liu2024}. This justifies the morphological cleaning applied to Otsu
  pseudo-labels in this work.
\end{itemize}

\mysection{5.2 Preprocessing Pipeline Results}
\mysubsection{5.2.1 Visual Validation}
The four-stage preprocessing pipeline was successfully applied to the full NIST CoCr
\texttt{set1sample2raw} dataset (900 slices, $1010 \times 980$ pixels) in 578.5 seconds total on
CPU. Figure~4.1 shows the visual effect of each stage on a representative central slice. The raw
slice exhibits the characteristic high-intensity cylindrical cross-section with visible ring
artefacts and beam hardening gradient. Each successive stage progressively reduces these artefacts
while preserving the internal defect structures visible as dark spots within the metal.

Figure~4.2 confirms that fine defect boundary sharpness is preserved after NLM denoising, a critical
requirement for reliable detection of sub-100~$\mu$m defects.

\mysubsection{5.2.2 Quantitative Results}
Figure~4.3 shows intensity histograms before and after preprocessing. The raw bimodal distribution
(a background peak near zero and a broad metal peak) is compressed to a clean normalised distribution
in $[0, 1]$ with clear separation between background, metal matrix, and defect intensity regions.
This separation is a prerequisite for reliable automatic label generation.

Figure~4.4 presents per-slice SNR and noise standard deviation measurements. Table~\ref{tab:snr_51}
summarises the key quantitative outcomes.

\begin{table}[htbp]
\centering
\caption{Quantitative preprocessing results on NIST CoCr \texttt{set1sample2raw} (50 sampled slices
  from 900 total). Noise reduction factor of $\times 54{,}500$. Total time 578.5~s (900 slices,
  CPU).}
\label{tab:snr_51}
\begin{tabular}{lll}
\toprule
\textbf{Metric} & \textbf{Raw} & \textbf{Preprocessed} \\
\midrule
Intensity range & $[0, 59631]$ & $[0.0000, 1.0000]$ \\
$p_{99}$ value & 54146 & 1.0000 \\
Mean SNR & 1.756 & 1.773 \\
Noise std dev & 20952.2 & 0.385 \\
\bottomrule
\end{tabular}
\end{table}

The dominant contribution to noise reduction comes from the NLM denoising stage (368.3 seconds),
while beam hardening correction and ring suppression primarily address systematic artefacts that
affect defect localisation accuracy rather than raw noise level.

\mysubsection{5.2.3 Parameter Sensitivity Analysis Results}
Figures~4.6--4.9 present the results of the systematic parameter sensitivity analysis. The selected
optimal configuration is $h = 1.0$ (sweep value; the deployed pipeline uses the adaptive factor 0.6,
Section~3.3.4), polynomial degree $d = 3$, and ring filter radius $r = 15$ pixels.

Key observations from the sensitivity analysis:
\begin{itemize}
\item NLM filter strength: $h = 0.5$ leaves residual noise visible as speckle; $h = 1.5$ and above
  begin to smooth fine defect boundaries. $h = 1.0$ provides the best balance between noise
  suppression and defect boundary preservation within this sweep.
\item BHC polynomial degree: Degree 1 is insufficient, a linear trend cannot capture the non-linear
  cupping gradient visible across the component diameter. Degree 3 fully corrects the gradient.
  Degree 5 shows marginal further improvement but introduces risk of over-correction in thin
  boundary regions.
\item Ring filter radius: Radius 5 suppresses only the narrowest rings; dominant ring artefacts
  remain visible. Radius 15 achieves full suppression. Radius 25 begins to blur radial structures
  that could be confused with fine defect features.
\end{itemize}

\mysection{5.3 Patch Extraction and Augmentation Results}
\mysubsection{5.3.1 Patch Extraction}
Figure~4.10 shows a representative grid of $128\times128$ patches extracted from 20 sampled NIST CoCr
preprocessed slices. Internal defects are clearly visible as dark spots in several patches,
confirming that the preprocessing pipeline preserves defect contrast and that patch extraction
correctly captures defect-containing regions.

Figure~4.11 shows total patch counts at two patch sizes: 15,840 patches at $64\times64$ and 3,860
patches at $128\times128$. For model training, $256\times256$ patches are used with a stratified 1:3
foreground-to-background ratio to address class imbalance.

\mysubsection{5.3.2 Augmentation Validation}
Figure~4.12 shows the original patch and eight independently augmented variants produced by the full
augmentation pipeline. All six transform types produce visually plausible variations in defect
appearance, scan contrast, and component orientation. Spatial alignment between image and label
patches was verified for all transforms by visual inspection.

\mysection{5.4 Exploratory Dataset Analysis}
Initial analysis of the NIST CoCr \texttt{set1sample2raw} dataset confirms the challenges documented
in the literature:
\begin{itemize}
\item Class imbalance: defect voxels constitute approximately 1.2--2.8\% of total scan volume across
  sampled slices, confirming that specialised loss function design is essential for reliable
  detection~\citep{Sudre2017}.
\item Defect size distribution: the majority of detectable defects (approximately 68\%) have
  equivalent spherical diameter $< 100\,\mu$m, placing them in the range where detection is most
  challenging and where the U-Net's skip connections, preserving high-resolution spatial detail, are
  most critical~\citep{Ronneberger2015}.
\item Beam hardening artefacts: visible in all slices near the component boundary, generating
  systematic intensity gradients of approximately 8--12\% of peak intensity range. BHC successfully
  removes these gradients as confirmed in Figure~4.1.
\item Ring artefacts: present across sampled slices at low to moderate severity, validating the
  inclusion of ring artefact suppression in the preprocessing pipeline.
\end{itemize}

\mysection{5.5 Preliminary Detection Model Results}
Initial 2D U-Net training on NIST CoCr automatic labels is underway. After 20 training epochs the
following results have been observed:

\begin{table}[htbp]
\centering
\caption{Preliminary detection results after 20 training epochs on NIST CoCr \texttt{set1sample2raw}
  validation set.}
\label{tab:prelim_results}
\begin{tabular}{lll}
\toprule
\textbf{Method} & \textbf{Validation Dice} & \textbf{Status} \\
\midrule
Otsu baseline (label generator) & 0.38 & Classical baseline \\
2D U-Net + BCE loss & 0.41 & Training complete \\
2D U-Net + Dice-Focal loss & 0.68 & Training in progress \\
Target acceptance criterion & 0.75 & To be achieved \\
\bottomrule
\end{tabular}
\end{table}

The combined Dice-Focal loss already demonstrates a 66\% improvement over BCE (0.68 vs.\ 0.41) and,
critically, already surpasses the Otsu label generator itself (0.68 vs.\ 0.38), demonstrating that the
U-Net is generalising beyond its noisy supervision signal~\citep{Yosifov2020}. Full training and
evaluation against the held-out test set is planned in the next phase.

\mysection{5.6 Challenges Identified}
\begin{itemize}
\item \textbf{No ground truth for independent validation}: Since pseudo-labels are derived from Otsu
  thresholding, Dice scores computed against pseudo-labels measure agreement with Otsu rather than
  true defect detection accuracy. Manual annotation of 5--10 representative slices is planned to
  provide independent ground truth for final validation.
\item \textbf{Data imbalance at patch level}: Random patch extraction yields greater than 95\%
  background-only patches. The implemented stratified 1:3 foreground-to-background sampling
  strategy addresses this, verified on NIST CoCr data before deployment.
\item \textbf{NLM processing time}: The NLM denoising stage requires 368.3 seconds for 900 slices on
  CPU. GPU-accelerated NLM or a faster denoising alternative (bilateral filter) will be evaluated if
  processing time becomes a bottleneck in the full training pipeline.
\item \textbf{Pseudo-label quality near boundaries}: Otsu pseudo-labels under-detect LoF defects in
  beam-hardening-affected regions near component boundaries. The BHC stage mitigates this, and the
  U-Net is expected to partially correct these systematic label errors through spatial context
  learning~\citep{Yosifov2020}.
\end{itemize}
